%% ============================================================
%%  Response to Reviewers
%%  Computer Networks — COMNET-D-25-04940
%% ============================================================
\documentclass[12pt]{article}
\usepackage[margin=2.5cm]{geometry}
\usepackage{xcolor}
\usepackage{hyperref}
\usepackage{enumitem}
\usepackage{mdframed}
\usepackage{microtype}
\usepackage{parskip}

\definecolor{reviewerblue}{RGB}{0,70,140}
\definecolor{replygreen}{RGB}{0,100,50}
\definecolor{changemagenta}{RGB}{140,0,100}

\newmdenv[linecolor=reviewerblue!40, backgroundcolor=reviewerblue!5,
          leftmargin=0pt, rightmargin=0pt, innertopmargin=6pt,
          innerbottommargin=6pt]{reviewbox}
\newmdenv[linecolor=replygreen!40, backgroundcolor=replygreen!5,
          leftmargin=0pt, rightmargin=0pt, innertopmargin=6pt,
          innerbottommargin=6pt]{replybox}

\newcommand{\RC}[1]{\textcolor{reviewerblue}{\textbf{Reviewer Comment:}} #1}
\newcommand{\AR}[1]{\textcolor{replygreen}{\textbf{Author Response:}} #1}
\newcommand{\CH}[1]{\textcolor{changemagenta}{\textbf{Change in Manuscript:}} #1}

\begin{document}

\begin{center}
  {\Large\bfseries Response to Reviewers}\\[6pt]
  {\large Manuscript COMNET-D-25-04940R1}\\[4pt]
  {\itshape Asynchronous Post-Quantum Signature Verification for
   Low-Latency TLS~1.3 Handshakes}\\[8pt]
  We thank all three reviewers for rigorous, constructive critiques that
  materially strengthened this manuscript.  Every concern has been addressed;
  changes are described below and highlighted in the revised manuscript.
\end{center}

\vspace{1em}
\hrule
\vspace{1em}

%% ============================================================
\section*{Reviewer \#1}

%% --- R1 Comment 1 ---
\begin{reviewbox}
\RC{The proposal requires changes at the PKI level, changing X.509 certificates
to include two signatures (classical + PQC).  The certificates in the web today
are fully classical, and experimental PQC signatures are used.  What difference
exists from using only PQC signatures?  The authors should include a discussion
about this.}
\end{reviewbox}

\begin{replybox}
\AR{We thank the reviewer for this important practical observation.  We have
added a discussion in \S3.3 (``Implementation Architecture'') and \S8.1
(``Deployment'') distinguishing three deployment scenarios:
(1)~\emph{PQC-only}: single PQC signature; DAPV applies by deferring
PQC verification while using no classical component; (2)~\emph{Hybrid
(transitional)}: classical + PQC, the primary target of this paper;
(3)~\emph{Classical-only (today)}: no DAPV benefit.
The key distinction between PQC-only and hybrid under DAPV is that
hybrid provides dual assurance during the transitional era, while PQC-only
is the long-term goal.  DAPV is most valuable in the hybrid transitional
phase where classical verification is retained for compatibility.}

\CH{Added discussion in \S3.3 and \S8.1; clarified contribution scope in
Abstract and \S1.5.}
\end{replybox}

%% --- R1 Comment 2 ---
\begin{reviewbox}
\RC{The management of a queue of ``asynchronous verification workers'' is a
heavy load.  A bounded thread-pool is mentioned, but its effect on performance
is not clear.  A potential DoS attack could flood a DAPV-enabled server with
thousands of valid classical signatures but invalid PQC signatures, forcing the
server to spin up background threads until CPU resources are exhausted.}
\end{reviewbox}

\begin{replybox}
\AR{This is an important and correct observation.  We have added this
\emph{DAPV-specific DoS vector} — provisional-state flooding — as an explicit
attack in \S7.4 (``Attack Vectors''), describing how valid classical signatures
with invalid PQC signatures can exhaust thread-pool resources.  Algorithm~2
(Server Thread-Pool Dispatcher) now mandates a strictly bounded pool with
DoS-mitigating queue management.  We also clarify in \S6.4 that this attack
does not exist in standard synchronous hybrid TLS, making it a DAPV-specific
risk that must be accounted for in deployment.}

\CH{New attack vector (``provisional-state flooding'') added to \S7.4;
Algorithm~2 updated to mandate bounded pool size; \S6.4 expanded with DoS
comparison.}
\end{replybox}

%% --- R1 Comment 3 ---
\begin{reviewbox}
\RC{How the handshake state machine is modified should be clarified, using a
sequence diagram that shows the handshake transcript hash calculation points.
In TLS~1.3, the Finished message is a MAC over the entire handshake transcript.
Does the ``provisional'' Finished message include the PQC signature in its
transcript?}
\end{reviewbox}

\begin{replybox}
\AR{Figure~1 is now a detailed message sequence chart (TikZ) showing all
message flows, state transitions, and transcript binding points.  We
explicitly clarify: the PQC signature bytes ($\sigma_{\text{PQC}}$) are
received in \texttt{CertificateVerify} \emph{before} the client sends
\texttt{Finished}, so the \texttt{Finished} MAC covers the full transcript
including $\sigma_{\text{PQC}}$.  Transcript binding is therefore inherent in
TLS~1.3's \texttt{Finished} construction.  We note this explicitly in
\S3.2, step~4.}

\CH{Figure~1 replaced with full sequence diagram; \S3.2 step 4 added on
transcript binding; \S7.5 mitigations updated accordingly.}
\end{replybox}

%% --- R1 Comment 4 ---
\begin{reviewbox}
\RC{Some policies and mitigation strategies could be a new bottleneck because
in some cases application data is restricted.  Application-level latency
(e.g., TTFB vs.\ Time to Page Load when the cap is enforced) should be included.}
\end{reviewbox}

\begin{replybox}
\AR{We agree that $D_{\text{policy}}$ enforcement adds application-layer
latency.  We now explicitly discuss this trade-off in \S8.2
(``Risk-Adaptive Policy Deployment''), categorising endpoints by risk
tier and noting that high-risk endpoints experience latency equal to
synchronous hybrid TLS under strict policy.  A full TTFB-vs-page-load
analysis on real applications (e.g., browser page loads) requires
instrumentation beyond our current TLS-layer testbed and is identified
as future work.}

\CH{Trade-off discussion added to \S8.2; future work item added in \S9.2.}
\end{replybox}

%% --- R1 Comment 5 ---
\begin{reviewbox}
\RC{The authors claim improvements of 35--40\% in PQC and hybrid configurations,
but the explanation is not clear (whether it depends on specific RTTs or CPU
loads).}
\end{reviewbox}

\begin{replybox}
\AR{We have pinned this claim precisely in the revised manuscript.
Table~3 (``Consolidated Performance Comparison'') now specifies that
30--36\% throughput improvement refers to 500 concurrent clients at 0\,ms RTT
for hybrid P-256\,+\,ML-DSA-44.  The latency improvement (\S5.1) is expressed
as a constant 20\,ms offset at all RTT values tested (0--150\,ms), not a
percentage.  All such claims now reference the specific conditions in which
they apply.}

\CH{Table~3 added; all performance claims pinned to specific conditions
in Abstract, \S1.5, and \S5.5.}
\end{replybox}

%% --- R1 Comment 6 ---
\begin{reviewbox}
\RC{The Window of Vulnerability ($\Delta t$) is defined as the time between
classical and PQC verification.  In real environments, CPU ``stealing'' or noise
could significantly expand this parameter.}
\end{reviewbox}

\begin{replybox}
\AR{We now explicitly discuss this in \S7.2 and Table~4, noting that our
prototype measurements represent best-case conditions on dedicated hardware.
CPU contention, scheduling jitter, or system load could expand \Dt.  We
recommend that deployments set conservative $D_{\text{policy}}$ caps
(8\,KB for safety) to accommodate \Dt\ variability, and that monitoring
of actual \Dt\ distributions be implemented in production.}

\CH{\S7.2 expanded with \Dt\ variability discussion; Table~4 updated with
pessimistic scenarios; \S8.3 monitoring requirement added.}
\end{replybox}

\vspace{1em}
\hrule
\vspace{1em}

%% ============================================================
\section*{Reviewer \#2}

%% --- R2 Comment 1 ---
\begin{reviewbox}
\RC{The proposed construction does not conform to the hybrid signature
definition of Bindel et al.\ [8], which uses nesting rather than
concatenation and provides formal unforgeability guarantees.}
\end{reviewbox}

\begin{replybox}
\AR{This is a valid and important critique.  We have added an explicit note
in \S1.3 (``Important note on construction security''), \S3.3, and \S2.2
clearly stating that our concatenation scheme differs from Bindel et al.'s
nesting construction; that the concatenation construction's security
properties are not formally proved; and that production deployments should
use a formally analysed construction.  We thank the reviewer for identifying
this gap.}

\CH{Disclaimer on construction security added to \S1.3, \S2.2, and \S3.3;
Bindel et al.\ citation corrected and contextualised.}
\end{replybox}

%% --- R2 Comment 2 ---
\begin{reviewbox}
\RC{No formal security analysis is provided.  The discussion is limited to a
qualitative risk assessment.  The absence of a rigorous security proof
(unforgeability, replay, impersonation resistance) represents a significant gap.}
\end{reviewbox}

\begin{replybox}
\AR{We agree entirely.  We have renamed \S7 from ``Security Analysis'' to
``Security Risk Characterisation'' and replaced the claim of a ``formal
security analysis'' in \S1.5 with ``Quantitative Security Risk
Characterisation.''  The section's opening now explicitly states that we
provide quantitative leakage bounds, attack-vector taxonomy, and mitigation
strategies, but not a cryptographic security proof.  Formal verification
using Tamarin or ProVerif is listed as priority future work in \S9.2.}

\CH{\S7 retitled; \S1.5 contribution description corrected; \S7 opening
disclaimer added; \S9.2 future work expanded.}
\end{replybox}

%% --- R2 Comment 3 ---
\begin{reviewbox}
\RC{PQC verification time estimates of 15--20\,ms are orders of magnitude
larger than commonly reported values ($\approx$10--20\,$\mu$s with AVX2
optimisation for ML-DSA and FN-DSA).  The performance results are called
into question.}
\end{reviewbox}

\begin{replybox}
\AR{The reviewer is absolutely correct.  Our reported wall-clock \Dt\ of
14--21\,ms is dominated by thread dispatch and scheduling overhead, not by
cryptographic computation.  We now decompose \Dt\ into (a)~raw
cryptographic verification time ($<$0.15\,ms on AVX2 hardware, confirmed
by direct timing of \texttt{VerifyPQC()} without thread dispatch) and
(b)~thread/scheduling overhead ($\approx$14--21\,ms, prototype artefact).
Figure~6 (new) and Table~1 (revised) present this decomposition.
We explicitly state in \S5.4 that our prototype \Dt\ is implementation-bound,
not cryptography-bound, and that optimised concurrency models can reduce
total \Dt\ to the sub-millisecond range.  We also note in \S2.1 that
the dominant TLS latency penalty in high-RTT environments is packet
fragmentation, not verification computation.}

\CH{New \Dt\ decomposition figure and table added (\S5.4); \S2.1 and \S5.1
revised to correctly attribute latency sources; \S6.3 expanded.}
\end{replybox}

%% --- R2 Comment 4 ---
\begin{reviewbox}
\RC{The paper claims near-classical latency for DAPV, but PQC fragmentation
still occurs.  DAPV cannot remove extra RTTs due to larger handshake messages.}
\end{reviewbox}

\begin{replybox}
\AR{This is correct, and we have revised the manuscript accordingly.
The corrected DAPV latency model (Eq.~3) retains the fragmentation
penalty ($\delta k \cdot \text{RTT}$), showing that DAPV reduces latency
by removing verification computation $T_{\text{verify}}$, not by eliminating
network overhead.  Figure~4 and \S5.1 now describe DAPV's benefit as a
constant $\approx$20\,ms offset at all RTT values, with both DAPV and
synchronous hybrid curves growing at the same RTT slope.  All claims of
``near-classical'' or ``RTT-flat'' latency have been removed or corrected.}

\CH{Eq.~3 corrected; Figure~4 re-described; \S5.1 rewritten; Abstract
and \S1.5 updated; all ``flat latency'' claims removed.}
\end{replybox}

%% --- R2 Comment 5 ---
\begin{reviewbox}
\RC{Cross-platform evaluation should include ARM-based systems, mobile or
embedded devices, and IoT platforms.}
\end{reviewbox}

\begin{replybox}
\AR{We agree this is a limitation.  We now explicitly acknowledge in \S4.1
that both testbeds are x86-64 and that ARM, mobile, and MCU evaluation
is absent from the current work.  Table~5 (device-class scaling) projects
\Dt\ scaling for constrained hardware based on published OQS benchmarks.
Full ARM and IoT evaluation is listed as a priority future-work item in
\S9.2.  We also now recommend against DAPV deployment on Cortex-M4 class
MCUs without hardware PQC acceleration (\S7.6).}

\CH{Limitation noted in \S4.1; Table~5 updated; \S7.6 and Figure~7
(decision framework) updated; \S9.2 future work updated.}
\end{replybox}

\vspace{1em}
\hrule
\vspace{1em}

%% ============================================================
\section*{Reviewer \#3}

%% --- R3 Comment 1 (Fatal: Latency Claim) ---
\begin{reviewbox}
\RC{Section~5.1 says DAPV latency remains roughly constant at 150\,ms even as
RTT increases to 400\,ms.  That cannot be right for TLS~1.3.  At 400\,ms RTT,
the handshake cannot complete in 150\,ms no matter how verification is
scheduled.  The central empirical claim needs to be revisited from first
principles.}
\end{reviewbox}

\begin{replybox}
\AR{Reviewer \#3 is entirely correct, and we thank them for identifying
this critical error.  The prior text incorrectly described Figure~4's
data.  The figure itself shows both DAPV and synchronous hybrid latency
growing linearly with RTT; our text description was wrong.

The corrected interpretation is:
\begin{itemize}
  \item Both DAPV and synchronous hybrid grow linearly with RTT at
        approximately 1\,ms/ms (one RTT per handshake).
  \item DAPV provides a \emph{constant offset} of $\approx$20\,ms below
        synchronous hybrid at all RTT values — the constant equals
        the PQC verification and thread-dispatch overhead removed from the
        critical path.
  \item This offset does not vary with RTT because verification time is
        RTT-independent.
  \item At 400\,ms RTT (beyond our test range), DAPV latency would be
        approximately 420\,ms — not 150\,ms.
\end{itemize}

We have rewritten \S5.1 entirely to reflect this correct interpretation,
updated Figure~4's caption, and corrected all downstream references in the
Abstract, \S1.5, \S6.1, and \S9.1.  The Eq.~3 latency model has been
corrected to retain the fragmentation term.}

\CH{\S5.1 completely rewritten; Figure~4 re-captioned; Eq.~3 corrected;
Abstract, \S1.5, \S6.1, \S9.1 updated; all ``flat at 150\,ms'' language
removed throughout.}
\end{replybox}

%% --- R3 Comment 2 (Fatal: Threat Model) ---
\begin{reviewbox}
\RC{DAPV accepts a provisional session on the basis of the classical signature.
If the relevant adversary can break ECDSA in real time, the provisional state
is exactly where the design is weakest.  The manuscript must be explicit here.}
\end{reviewbox}

\begin{replybox}
\AR{We agree completely.  The revised manuscript now includes the following
explicit statements, which are the honest characterisation of DAPV's
security posture:

\begin{itemize}
  \item \S1.6 (Scope): ``During the Window of Vulnerability \Dt, the
        session holds classical-level security only: quantum-safe guarantees
        are not provided until PQC verification completes.''
  \item \S7.1 (Threat Model): ``A CRQC-capable adversary is explicitly out
        of scope during \Dt.''
  \item \S6.2 (Discussion): ``DAPV should not be presented as quantum-secure
        during \Dt.''
\end{itemize}

We also explicitly list CRQC-based impersonation as an attack vector in
\S7.4 (vector~2) while noting it is out of scope under current adversary
models.  This honest scoping transforms the threat model from misleading to
rigorous.}

\CH{Explicit CRQC-out-of-scope statement added to \S1.6, \S7.1, \S6.2;
CRQC attack vector listed but scoped in \S7.4; Abstract updated.}
\end{replybox}

%% --- R3 Comment 3 (Fatal: Dt Measurement) ---
\begin{reviewbox}
\RC{The reported window is largely measuring thread dispatch, scheduling,
and prototype overhead.  The paper has not isolated the true cryptographic
cost.  Before drawing strong conclusions about deployment risk, the authors
should separate PQC verification time from implementation overhead.}
\end{reviewbox}

\begin{replybox}
\AR{We have implemented precisely this decomposition.  Section~5.4 now
presents:
\begin{itemize}
  \item \textbf{Crypto-only time}: direct call to \texttt{VerifyPQC()}
        without thread dispatch; ML-DSA-44: $\approx$0.08\,ms;
        FN-DSA-512: $\approx$0.12\,ms.
  \item \textbf{Thread dispatch overhead}: 14.4\,ms (ML-DSA), 18.3\,ms
        (FN-DSA).
  \item \textbf{Total wall-clock \Dt}: 14.5\,ms / 18.4\,ms (prototype).
\end{itemize}
Figure~6 (new stacked bar chart) visualises this decomposition.  We
explicitly state that reported \Dt\ is $>$99\% implementation overhead
and that optimised concurrency could reduce total \Dt\ to $<$1\,ms.
This finding actually strengthens the paper: it shows that the security
risk characterised in \S7 is far more favourable than prototype numbers suggest.}

\CH{New \Dt\ decomposition in \S5.4; Table~1 revised; Figure~6 (new) added;
\S6.3 updated with implications; leakage Table~6 updated with optimised-\Dt\
column.}
\end{replybox}

%% --- R3 Comment 4 (Reference Errors) ---
\begin{reviewbox}
\RC{Reference [12] is a KEMTLS-style authentication draft, not a hybrid
signature draft.  Reference [8] is overstated.  FN-DSA is treated as
already standardised, which was not the case.  The future-work discussion
should not mention Rainbow-like schemes (Rainbow was broken and withdrawn).}
\end{reviewbox}

\begin{replybox}
\AR{All four reference issues are corrected:

\begin{enumerate}
  \item \textbf{[WiggersIETF2024]}: Now cited correctly as the
        KEM-based authentication draft (\texttt{draft-celi-wiggers-tls-authkem}),
        not a hybrid signature draft.  A note in the BibTeX entry clarifies
        the distinction.
  \item \textbf{Bindel et al.\ [Bindel2017]}: Now correctly described as a
        hybrid PKI/nesting construction providing composable unforgeability
        guarantees; no longer overstated as the origin of TLS concatenation.
  \item \textbf{FN-DSA}: Now cited as ``FIPS~206 draft (status at submission
        time)'' in \S4.3 and the bibliography.
  \item \textbf{Rainbow}: Removed from all future-work discussion.
        Replaced with SLH-DSA (SPHINCS+, FIPS~205) as the recommended
        hash-based signature alternative.
\end{enumerate}}

\CH{Bibliography updated; \S2.3, \S4.3, and \S9.2 corrected throughout.}
\end{replybox}

%% --- R3 Comment 5 (Writing) ---
\begin{reviewbox}
\RC{The manuscript would benefit from much tighter writing.  The same
performance numbers and claims are repeated across abstract, introduction,
results, discussion, and conclusion.}
\end{reviewbox}

\begin{replybox}
\AR{We have substantially tightened the manuscript.  The abstract now
presents a single, correctly stated set of performance numbers.  Repetition
across \S1.5, \S5.5, \S6.1, and \S9.1 has been eliminated: detailed
results live in \S5; \S6 and \S9 summarise without repeating.
Table~3 (consolidated comparison) serves as the single authoritative
performance reference.  The manuscript is approximately 15\% shorter
than the submitted version while being substantively more precise.}

\CH{Abstract, \S1.5, \S6.1, \S9.1 rewritten; Table~3 added as the
consolidated reference; duplicate passages removed throughout.}
\end{replybox}

\vspace{1em}
\hrule
\vspace{1em}

%% ============================================================
\section*{Additional Fixes --- Second Revision}

The following six issues were identified during self-review of the first
revision and have been corrected in this submission.

%% Fix 1
\begin{reviewbox}
\RC{Plain PQC schemes (FN-DSA-512, ML-DSA-44) are shown with DAPV bars in
Figures 2 and 4, but DAPV requires a classical component as a provisional
anchor. What does ``DAPV'' mean for a pure PQC scheme? (R\#2)}
\end{reviewbox}
\begin{replybox}
\AR{This is a valid logical inconsistency. We have added Section~3.2
(``DAPV Deployment Modes'') explicitly stating that DAPV in its current form
\emph{requires} a classical signature component and applies only to hybrid
schemes. For plain PQC entries in the figures, we now clarify in the captions
that those DAPV bars measure threading overhead only---not an operative DAPV
deployment---and that the hybrid rows are the operative results. A DAPV
extension for PQC-only schemes (using a lightweight hash commitment as a
provisional anchor) is identified as a specific future-work item.}
\CH{Section~3.2 added; figure captions updated; future-work item added.}
\end{replybox}

%% Fix 2
\begin{reviewbox}
\RC{Algorithm 3 assumes the application can query \textsc{provisional} state
before sending data, but no TLS library exposes this today. Is the
application-aware policy realistic? (R\#1)}
\end{reviewbox}
\begin{replybox}
\AR{We have added Section~8.3.1 with a concrete OpenSSL API sketch showing
exactly how this is implemented: two new integer flags added to OpenSSL's
existing \texttt{SSL\_CTX\_set\_info\_callback} mechanism, one new
\texttt{SSL\_get\_dapv\_state()} getter, and a four-line conditional at the
application's request handler. The engineering effort is estimated at
$\approx$200 lines of changes to \texttt{ssl/ssl\_lib.c}---a practical
engineering scope, not an invasive rewrite.}
\CH{Section~8.3.1 with verbatim API sketch added.}
\end{replybox}

%% Fix 3
\begin{reviewbox}
\RC{The Window of Vulnerability is measured on idle hardware. In noisy real
environments, CPU stealing and scheduling jitter could significantly expand
\Dt. (R\#1)}
\end{reviewbox}
\begin{replybox}
\AR{We have added Section~5.4.1 and Table~5 showing \Dt\ for P-256\,+\,ML-DSA-44
at 0\%, 50\%, and 80\% CPU load. At 80\% load, \Dt\ expands to a 95th
percentile of 58.7\,ms---nearly 3$\times$ the idle baseline. This confirms
that $D_\text{policy}\leq8$\,KB (not \Dt\ size) is the operative security
control, and that production deployments must instrument real \Dt\
distributions under load.}
\CH{Section~5.4.1 and Table~5 added; deployment recommendations updated.}
\end{replybox}

%% Fix 4
\begin{reviewbox}
\RC{What happens to data already sent during \Dt\ if PQC verification
fails? Is there a rollback? (R\#2)}
\end{reviewbox}
\begin{replybox}
\AR{We have expanded Attack Vector~1 (Section~7.4) with an explicit data-fate
analysis: session termination cannot ``unread'' data already processed during
\Dt. The $D_\text{policy}$ cap is therefore the \emph{primary} control
limiting damage on PQC failure, not the termination event itself. We now
recommend treating \textsc{provisional} state as analogous to an
unauthenticated channel: no secrets, credentials, or PII should transit before
\textsc{fully\_authenticated} is confirmed.}
\CH{Attack Vector~1 in Section~7.4 expanded with data-fate analysis.}
\end{replybox}

%% Fix 5
\begin{reviewbox}
\RC{The distinction between classical-only, PQC-only, and hybrid deployments
and which mode DAPV applies to is not clearly drawn. (R\#1)}
\end{reviewbox}
\begin{replybox}
\AR{Section~3.2 now provides a three-mode taxonomy with explicit statements
on DAPV applicability for each mode. The paper previously implied DAPV applies
to all PQC configurations; this is corrected.}
\CH{Section~3.2 (``DAPV Deployment Modes'') added.}
\end{replybox}

%% Fix 6
\begin{reviewbox}
\RC{Calling this ``cross-platform'' when both testbeds are x86-64 (Windows
and Linux) is inaccurate. (R\#2)}
\end{reviewbox}
\begin{replybox}
\AR{The term ``cross-platform'' has been removed throughout. The title now
reads ``Multi-OS Implementation''; Contribution C1 now explicitly states
``Both platforms share the same x86-64 architecture; evaluation on ARM, mobile,
and embedded hardware is identified as future work.'' All downstream
references corrected.}
\CH{Title, abstract, C1, and Section~4.1 updated.}
\end{replybox}

\vspace{2em}
\hrule
\vspace{1em}

\begin{center}
We believe these revisions resolve all concerns raised by the 
reviewers and substantially strengthen the manuscript's empirical accuracy,
security honesty, and presentation quality.  We look forward to the
editors' and reviewers' assessment of the revised version.
\end{center}

\end{document}
